\documentclass[12pt,draftclsnofoot,onecolumn]{IEEEtran}

\usepackage{amsmath,amssymb}
\usepackage{booktabs}
\usepackage{graphicx}
\usepackage{multirow}
\usepackage{threeparttable}
\usepackage{longtable}
\usepackage{array}
\usepackage{tabularx}
\usepackage{siunitx}
\usepackage{xcolor}
\usepackage[hidelinks]{hyperref}
\usepackage{microtype}

\newcolumntype{L}[1]{>{\raggedright\arraybackslash}p{#1}}

\sisetup{detect-all,round-mode=places,round-precision=3}
\graphicspath{{figures/}}

\title{Reliability of COVID-19 Respiratory-Audio Screening Under Temporal, Metadata, and Cross-Dataset Shift}

% Verify spelling, order, ORCID identifiers, and affiliations before circulation.
\author{Nishant Harkut and Ritu Tiwari}

\begin{document}
\maketitle

\begin{abstract}
Respiratory audio has been proposed as a low-cost input for coronavirus disease
2019 (COVID-19) screening, although reported performance is sensitive to cohort
construction and evaluation design. We studied cough, breathing, and speech
recordings from Coswara under participant-disjoint internal, calendar-aware,
chronological, and cross-dataset protocols. External testing used 8,331
independently collected COUGHVID cough recordings. The selected internal
multimodal system reached an area under the receiver operating characteristic
curve (AUROC) of 0.897 (95\% bootstrap confidence interval 0.854--0.935). In
modality-matched transfer, four conventional cough classifiers lost
approximately 0.30 AUROC. WavLM changed from 0.812 internally to 0.484
externally, and a
spectrogram CNN--BiGRU changed from 0.737 to 0.548. Only 110 of the top 800
features selected in early and late periods overlapped (Jaccard index 0.074).
Within Coswara, metadata alone reached AUROC 0.964, whereas full retraining
after label permutation produced mean AUROC 0.503. In an exploratory aligned
analysis ($n=61$), adding audio to the full metadata model changed AUROC by
0.012 (95\% confidence interval $-0.033$ to 0.053; paired DeLong $p=0.572$).
Strong internal discrimination therefore did not establish temporal stability,
incremental value, or cross-dataset portability. This retrospective reliability
study does not establish a deployable diagnostic system.
\end{abstract}

\begin{IEEEkeywords}
COVID-19, respiratory audio, multimodal learning, temporal validation,
external validation, dataset shift, calibration.
\end{IEEEkeywords}

\section{Introduction}

Smartphone recordings of cough, breathing, and speech offer a convenient way
to study respiratory health outside a clinic. Public datasets and early model
studies therefore motivated audio-based COVID-19 screening as a possible
low-cost triage aid \cite{xia2021covid,bhattacharya2023coswara,orlandic2021coughvid}.
The central methodological difficulty is that these datasets were assembled
during a rapidly changing pandemic through heterogeneous recruitment,
self-report, devices, locations, and testing practices. A classifier can
separate labels in one sample without learning a stable disease-related audio
pattern that transfers to another period or collection system.

Several strands of prior work illustrate the tension. Architecture-led studies
reported strong within-dataset discrimination using transfer learning,
multimodal score fusion, and hierarchical spectrogram transformers
\cite{pahar2022transfer,chetupalli2023multimodal,aytekin2024hst}. In contrast, realistic
participant-independent evaluation showed that biased split construction can
raise AUROC from approximately 0.71 to 0.90 \cite{han2022realistic}. Coppock
\emph{et al.} found that an audio model's AUROC decreased from 0.846 to 0.619
after matching measured confounders, and that symptom models were strong
comparators \cite{coppock2024audio}. Temporal drift and cross-dataset studies
have since reinforced that within-dataset performance and deployment-time
transportability are different estimands \cite{ganitidis2025drift,islam2026robust}.

This study asks a narrower question than whether one more classifier can raise
an internal benchmark. It asks which conclusions survive when models developed
from the same source data are examined under different sources of shift:
participant allocation, calendar time, label prevalence, recording protocol,
and dataset origin. We use Coswara because it contains multiple sound types per
participant, and COUGHVID as a separately collected cough-only target. We also
test whether non-audio variables predict the label, whether audio adds value to
those variables on aligned participants, and whether calibration or target
threshold adjustment can repair weak external discrimination.

The contributions are:

\begin{enumerate}
    \item a participant-audited multimodal pipeline spanning conventional
    acoustic functionals and learned audio representations;
    \item protocol-specific internal, calendar-aware, chronological, and
    modality-matched external evaluations with uncertainty estimates;
    \item negative controls, metadata prediction, feature-stability, and
    domain-overlap analyses that distinguish observed associations from
    stronger causal interpretations; and
    \item base-rate-aware calibration, fixed-sensitivity, and decision analyses
    that separate ranking performance from operational usefulness.
\end{enumerate}

The study is not a clinical validation and does not claim that COVID-19 leaves
no acoustic signature. Its purpose is to determine how much evidence the
available retrospective benchmarks provide for a reliable screening model.

\section{Related Work}

\subsection{Respiratory-Audio Datasets and Benchmark Models}

Coswara contains multiple respiratory and vocal tasks with participant
metadata. COUGHVID is a large, heterogeneous cough corpus
\cite{bhattacharya2023coswara,orlandic2021coughvid}. The COVID-19 Sounds resource
similarly includes breathing, cough, and voice and introduced realistic
participant-independent baselines \cite{xia2021covid}. These resources differ
in recruitment, label construction, class prevalence, sound prompts, and
quality control. Dataset identity is therefore part of the experimental
condition, not an interchangeable filename.

Early systems used MFCCs, spectral descriptors, OpenSMILE functionals, support
vector machines, convolutional networks, and score fusion. Feature-ranking
work showed that the useful representation can depend on the dataset
\cite{meister2022ranking}, and the DiCOVA study of Avila \emph{et al.} combined
6,373-dimensional ComParE features with a log-mel CNN
\cite{avila2021feature}. Multimodal Coswara work reported gains from averaging
predictions across respiratory and vocal tasks, but selected combinations from
many possible subsets \cite{chetupalli2023multimodal}. These studies motivate multimodal
modeling while also making validation-only fusion selection essential.

\subsection{Deep and Transformer Representations}

Transfer learning with CNN, LSTM, and ResNet representations has been evaluated
for cough, breathing, and speech \cite{pahar2022transfer}. WavLM provides a
self-supervised transformer representation learned from large speech corpora
\cite{chen2022wavlm}. Aytekin \emph{et al.} proposed the hierarchical
spectrogram transformer (HST), which partitions a $224\times224$ log-mel
spectrogram into patches and expands the attended context through local-window
attention and patch merging \cite{aytekin2024hst}. Their reported internal
results are important architecture references, but their Cambridge and
COUGHVID experiments use dataset-specific tasks and 10-fold partitions. They
are not direct estimates of Coswara-to-COUGHVID transfer.

The ESWA DNDT/DNDF study combines 193 cough features, recursive feature
elimination with cross-validation, Bayesian hyperparameter optimization,
SMOTE, and threshold moving \cite{islam2026robust}. It reports high
within-dataset AUCs for Coswara and COUGHVID, but its explicit cross-dataset
table also contains much lower source-to-target values. This distinction is
central to the present comparison: within-dataset cross-validation and frozen
external transfer answer different questions.

\subsection{Bias, Drift, and Clinical Utility}

Han \emph{et al.} demonstrated that participant overlap and demographic
imbalance can produce optimistic audio estimates \cite{han2022realistic}.
Coppock \emph{et al.} extended this analysis through prespecified matched and
longitudinal cohorts, symptom baselines, incremental models, and utility
analysis \cite{coppock2024audio}. Ganitidis \emph{et al.} studied model drift
and adaptation across development and later cough data \cite{ganitidis2025drift}.
Together, these studies shift the research question from maximum internal AUC
to transportability under a stated target population and time.

Prediction-model reporting guidance reaches the same conclusion. TRIPOD+AI
requires explicit data provenance, participant flow, modeling choices,
performance measures, and limitations \cite{collins2024tripod}; PROBAST+AI
emphasizes risk of bias and applicability across participants, predictors,
outcomes, and analysis \cite{moons2025probast}. We therefore report AUROC with
AUPRC, calibration, threshold-dependent metrics, and uncertainty, and keep
confirmatory comparisons separate from exploratory mechanism analyses.

\section{Materials and Methods}

\subsection{Study Design and Datasets}

The development source was Coswara. The indexed release contained 2,746
participants; 2,114 participants with resolved positive or negative labels
contributed 19,024 recordings to supervised analyses. The remaining 632
participants (5,688 recordings) were retained in data audits but excluded from
model fitting and evaluation. Each participant may contribute several
prompted recordings, including heavy and shallow cough, deep and shallow
breathing, counting, and sustained vowels. All recordings from a participant
were assigned to one partition. Rows with unresolved labels remained in the
data-quality audit but were excluded from supervised evaluation.

The external target contained 8,331 COUGHVID recordings: 285 positive and
8,046 negative under the processed \texttt{status\_SSL} analytic label, for a
known-label prevalence of 3.42\%. Each external recording had a unique
identifier; COUGHVID does not provide reliable participant linkage for this
cohort, so the external analysis unit is a recording rather than a participant.
No
COUGHVID sample was used for source feature ranking, model fitting, fusion
selection, or source threshold selection. Because COUGHVID contains coughs,
external evaluation was restricted to the Coswara cough branch.

The two datasets do not implement identical outcome ascertainment. Coswara's
analytic labels originate from participant COVID-19 status fields, whereas the
COUGHVID target mapping uses the processed \texttt{status\_SSL} field after the
project's predefined quality and label rules. The raw COUGHVID \texttt{status}
field was available for fewer recordings and disagreed with \texttt{status\_SSL}
for 175 recordings with both fields resolved; a raw-status sensitivity analysis
remains outstanding. Label-definition mismatch is therefore a plausible
contributor to transfer error.

\subsection{Audio Quality and Preprocessing}

Audio paths, participant identifiers, modalities, labels, durations, and
quality flags were indexed before modeling. Loading converted supported source
formats to mono waveforms at 16~kHz and recorded decoding failures. Quality
checks identified 23,539 valid, 597 mostly silent, 390 corrupt, and 186 short
recordings in the indexed corpus. Conventional feature extraction removed
leading and trailing low-energy regions at a 35-dB threshold, localized an
active region by frame-level root-mean-square energy, peak-normalized the
signal, and center-cropped or padded cough, breathing, and speech to 4, 5, and
5~s, respectively. Time--frequency calculations used a 400-sample window,
160-sample hop, and 64 mel bands from 20 to 8,000~Hz. Each experiment declared
whether it used all decodable samples or only quality-passing samples; the rule
was fixed before feature computation and applied unchanged outside training.

\subsection{Acoustic Representations}

The main conventional representation concatenated three sources of
recording-level information: the existing acoustic bank, the ComParE 2016
functionals, and the INTERSPEECH 2010 Paralinguistic Challenge functionals.
The latter two banks were extracted with openSMILE
\cite{eyben2010opensmile,schuller2010is10,schuller2013compare}.
The resulting merged table contained 10,140 numeric candidates: 6,377
ComParE-derived columns, 1,586 IS10-derived columns, and 2,177 project acoustic
columns. LightGBM gain importance \cite{ke2017lightgbm} was calculated separately in each available
modality on the original source-training partition, normalized within modality,
and averaged across modalities. The primary reduced bank retained the 800
highest-ranked columns; the same names and order were frozen for validation,
test, chronological stress testing, and COUGHVID transfer. This design prevents
target-label feature selection but means the chronological analysis is not a
fully nested re-selection experiment.

Two learned representations tested whether the portability result was specific
to conventional functionals. WavLM Base Plus encoded short cough segments and
pooled segment predictions to the participant or recording level. The
spectrogram branch used a convolutional front end followed by a bidirectional
gated recurrent unit (CNN--BiGRU), so local time--frequency patterns were
learned before temporal aggregation. Both branches used source validation for
checkpoint and threshold decisions and were evaluated on COUGHVID without
target-label retraining.

\subsection{Classifiers and Imbalance Handling}

The conventional bank evaluated regularized radial-basis SVC, LightGBM,
XGBoost, and CatBoost classifiers
\cite{cortes1995svm,ke2017lightgbm,chen2016xgboost,prokhorenkova2018catboost}.
Imbalance handling was restricted to the
training partition. For model configurations denoted by an SMOTE fraction,
synthetic minority examples were generated only after the training subset and
feature space had been fixed; validation and test distributions were not
resampled \cite{chawla2002smote}. Hyperparameters and feature count were selected from training and
validation evidence, never from test or external labels.

\subsection{Multimodal Fusion}

Each modality model produced a probability $p_m(x)$ for participant $x$. Three
fusion rules were considered. Uniform fusion used
\begin{equation}
    p_{\mathrm{mean}}(x)=\frac{1}{|\mathcal{M}_x|}
    \sum_{m\in\mathcal{M}_x}p_m(x),
\end{equation}
where $\mathcal{M}_x$ is the set of available modalities. Validation-weighted
fusion transformed each modality's validation AUPRC as
$r_m=\max(\mathrm{AUPRC}_m-0.5,0.01)$ and normalized it,
\begin{equation}
    w_m=\frac{r_m}{\sum_j r_j},\qquad
    p_{\mathrm{weighted}}(x)=\sum_m w_m p_m(x).
\end{equation}
Stacked fusion fitted class-balanced L2 logistic regression only on aligned
validation predictions,
\begin{equation}
    p_{\mathrm{stack}}(x)
    =\sigma\!\left(\beta_0+\sum_m\beta_m p_m(x)\right),
\end{equation}
where $\sigma$ is the logistic function. The implementation consumes modality
probabilities directly; it does not first apply a logit transform. Fitted
coefficients and the validation-selected balanced-accuracy threshold were then
frozen for test evaluation. Consequently, stack coefficients were learned,
not hard-coded, and test labels did not determine them.

\subsection{Evaluation Protocols}

Four protocols were distinguished.

\begin{enumerate}
    \item \emph{Participant-disjoint internal evaluation} separated source
    participants across train, validation, and test partitions.
    \item \emph{Time-stratified participant evaluation} retained participant
    separation while distributing calendar periods across partitions. This
    tests a more calendar-balanced internal scenario.
    \item \emph{Early-to-late evaluation} trained on an earlier calendar block,
    selected models on source validation data, and evaluated later
    participants. The frozen development feature bank described above was used;
    feature selection was not refitted solely inside the early block. A reverse
    late-to-early analysis was exploratory.
    \item \emph{External transfer} trained and selected the cough branch on
    Coswara and applied it unchanged to COUGHVID.
\end{enumerate}

The protocol-specific selected models are not automatically paired estimands.
The four-number descriptive ladder is therefore supplemented by
model-and-modality-matched cough comparisons and paired tests only where the
same participants have predictions from both models.

Figure~\ref{fig:study-design} summarizes the development and evaluation
boundaries. In particular, the chronological and external analyses reuse the
source-developed representation and source-validation decisions; they do not
select a new model from the later or external labels.

\begin{figure}[!t]
\centering
\includegraphics[width=0.99\linewidth]{fig01_study_design.pdf}
\caption{Study design. Coswara supplied multimodal source development data;
COUGHVID supplied a cough-only external target. Participant separation,
training-only resampling and ranking, and source-validation model and threshold
selection were maintained throughout.}
\label{fig:study-design}
\end{figure}

\subsection{Confounding and Shift Analyses}

Metadata-only logistic models used three predefined input sets: symptoms,
demographic and recording-protocol variables, and their safe combined set.
Permutation importance grouped encoded columns by their source concept.
Full retraining with permuted labels tested whether the observed metadata
performance could arise from software leakage.

A source-versus-target classifier and its predicted-probability band measured
dataset separability and limited overlap. The reported fraction
outside the source probability band is a diagnostic, not a formal proof of a
causal positivity violation.

Temporal feature stability was measured by the Jaccard index between early and
late top-800 sets,
\begin{equation}
    J(F_E,F_L)=\frac{|F_E\cap F_L|}{|F_E\cup F_L|}.
\end{equation}

\subsection{Performance and Statistical Analysis}

Discrimination was summarized by AUROC and AUPRC. Threshold-dependent results
included sensitivity, specificity, precision, balanced accuracy, and F1. The
threshold was chosen on validation data by balanced accuracy unless a stated
fixed-sensitivity analysis was performed. Calibration was summarized by Brier
score, negative log likelihood, and expected calibration error (ECE):
\begin{align}
    \mathrm{Brier} &= \frac{1}{n}\sum_{i=1}^{n}(p_i-y_i)^2,\\
    \mathrm{ECE} &= \sum_{b=1}^{B}\frac{|I_b|}{n}
    \left|\operatorname{acc}(I_b)-\operatorname{conf}(I_b)\right|.
\end{align}
The reliability-bin formulation follows established probability-calibration
analysis \cite{niculescu2005calibration}.

For the target-only recalibration sensitivity analysis, COUGHVID was divided
before fitting into 2,082 calibration recordings and 6,249 disjoint evaluation
recordings. Platt scaling and isotonic regression were fitted only on the
calibration subset and assessed only on the evaluation subset. This analysis
tested whether probability rescaling could improve calibration and
threshold-dependent performance; it was not a new external test because target
labels were used to fit the recalibrator.

Bootstrap intervals resampled the declared analysis unit. For external AUROC
differences, source participants and target recordings were sampled
independently; repeated Coswara cough recordings were kept within participant
clusters. DeLong comparisons were used only for paired
predictions on the same labeled participants \cite{delong1988auc}. Decision-curve net benefit at
threshold probability $p_t$ was
\begin{equation}
    \mathrm{NB}(p_t)=\frac{TP}{n}-\frac{FP}{n}\frac{p_t}{1-p_t}.
\end{equation}
This definition compares the model with treat-all and treat-none strategies
across threshold probabilities \cite{vickers2006dca}.
The selected internal system, matched cough transfer, and representation-family
transfer formed the primary reliability analyses in this report. Subgroups,
context controls, reverse temporal testing, and mechanism analyses were
post-development exploratory analyses; their $p$-values were not interpreted
as a family of confirmatory tests.

\section{Results}

\subsection{Internal and Calendar-Aware Evaluation}

The selected participant-disjoint multimodal model combined cough and speech
through a validation-fitted logistic stack. It achieved AUROC 0.897, AUPRC
0.863, balanced accuracy 0.825, and F1 0.752 on 314 test participants. Under the
time-stratified participant protocol, the selected uniform cough--breath--speech
fusion achieved AUROC 0.849, AUPRC 0.783, balanced accuracy 0.783, and F1 0.705
on 431 participants. These two rows describe protocol-specific selected models;
their numerical difference is not a paired estimate of the effect of calendar
stratification.

Across available seed replications, the validation-stacked internal
cough--speech fusion had mean AUROC $0.895\pm0.003$ (four seeds). The selected
early-to-late breath ensemble had mean AUROC $0.691\pm0.007$ (three seeds).
The early-to-late selected row reached AUROC 0.698 but had ECE 0.711, indicating
that the probabilities were poorly aligned with observed outcomes. Again,
these rows involve different final modality combinations and are evidence of
protocol sensitivity across the selected workflow, not a controlled paired
model comparison.

The exploratory reverse temporal run produced AUROC 0.920 but AUPRC 0.029,
F1 0.011, and ECE 0.471. The high AUROC therefore did not imply a useful
operating point under the reversed prevalence and cohort composition.

\subsection{Modality-Matched External Transfer}

Table~\ref{tab:external} reports like-for-like cough transfer. Every
conventional model lost substantial discrimination when moved from Coswara to
COUGHVID. In the independent two-sample bootstrap, Coswara cough predictions
were first averaged within each participant and participants were resampled;
COUGHVID recordings were resampled independently. All four 95\% intervals for
the AUROC difference excluded zero.

\begin{table}[!t]
\centering
\caption{Cough-only internal-to-external transfer. Conventional-model
differences include participant-cluster-correct 95\% bootstrap intervals.}
\label{tab:external}
\begin{threeparttable}
\begin{tabular}{lccc}
\toprule
Model family & Internal AUROC & External AUROC & Difference (95\% CI) \\
\midrule
LightGBM & 0.853 & 0.543 & 0.310 (0.247, 0.366) \\
SVC--RBF & 0.868 & 0.523 & 0.345 (0.291, 0.398) \\
CatBoost & 0.849 & 0.531 & 0.318 (0.256, 0.376) \\
XGBoost & 0.850 & 0.532 & 0.318 (0.262, 0.371) \\
WavLM Base Plus & 0.812 & 0.484 & 0.328 (not estimated) \\
CNN--BiGRU & 0.737 & 0.548 & 0.189 (not estimated) \\
\bottomrule
\end{tabular}
\begin{tablenotes}\footnotesize
\item The external prevalence was 3.42\% (285/8,331). COUGHVID evaluates only
the cough branch. Conventional internal values are participant-level means over
available Coswara cough recordings ($n=316$); external values use 8,331 target
recordings. Deep-model differences are descriptive because comparable
cluster-correct intervals were not estimated for those runs.
\end{tablenotes}
\end{threeparttable}
\end{table}

At required sensitivity of at least 0.90, the conventional external models
retained only approximately 0.11--0.16 specificity and 0.035--0.037 precision.
That precision was close to the 0.034 target prevalence. This means that, at
the chosen screening constraint, a positive prediction provided little gain
over the target base rate.

Target-only Platt and isotonic recalibration changed probability scale and
threshold-dependent metrics but did not materially change AUROC. The result
separates calibration error from ranking failure: threshold adjustment can
alter decisions, but cannot recover discrimination that is absent in the score
ordering.

Figure~\ref{fig:validation-transfer} keeps the protocol-specific selected
systems separate from the modality-matched transfer comparison. This prevents
the 0.897 multimodal internal value from being presented as the source endpoint
of a cough-only external experiment.

\begin{figure}[!t]
\centering
\includegraphics[width=0.99\linewidth]{fig02_validation_and_transfer.pdf}
\caption{Discrimination under the evaluated protocols. (A) Protocol-specific
selected systems; modality combinations differ and the points are descriptive.
(B) Cough-only internal-to-COUGHVID transfer across conventional and learned
representations. Brackets give 95\% bootstrap intervals for conventional-model
AUROC differences.}
\label{fig:validation-transfer}
\end{figure}

\subsection{Metadata Association and Incremental Audio Value}

The full safe metadata model reached AUROC 0.964; symptoms alone reached 0.932,
and demographic plus recording-protocol variables reached 0.914. After full
model retraining with permuted labels, mean AUROC values were 0.503, 0.503, and
0.503, respectively. The negative control argues against a simple leakage or
metric-calculation bug in the metadata pipeline.

Grouped coefficient magnitude and permutation analyses identified recording
year and recording-protocol variables as influential. This demonstrates that
label-predictive collection context exists. It does not by itself prove that
the audio model encoded the same variables, because the metadata and audio
predictors are distinct measurements.

Incremental value was evaluated on candidates selected by validation
performance before test comparison. For the highest-ranked aligned
cough--speech candidate, the full metadata model achieved test AUROC 0.924 and
metadata plus audio achieved 0.935. The difference was 0.012 (95\% bootstrap
interval $-0.033$ to 0.053; paired DeLong $p=0.572$; $n=61$). Against the
symptoms-only model, adding the same audio candidate changed AUROC by 0.063
(95\% bootstrap interval $-0.005$ to 0.149; paired DeLong $p=0.104$). These
small aligned samples do not establish absence of incremental audio value, but
they do not provide precise evidence of an improvement.

\subsection{Mechanism and Robustness Analyses}

Only 110 of the top 800 acoustic features ranked independently in the early
period were also ranked in the late period. The union contained 1,490 features,
giving a Jaccard overlap of 0.074. This sensitivity analysis indicates that the
feature-ranking output was temporally unstable. Because prevalence and case mix
also changed, it neither isolates the cause nor proves that underlying biology
changed.

A source-versus-target classifier based on 500 common features achieved AUROC
0.750. Using the 5th--95th percentile band of source-domain probabilities,
25.2\% of external recordings fell outside the source band. This diagnostic
supports meaningful source-target mismatch and warns that reweighting or
covariance alignment may require extrapolation.

Audio label-shuffle retraining produced AUROC values of 0.537--0.551 across the
three final protocols. These are close enough to chance to support the absence
of gross label leakage, while their finite-sample deviation from exactly 0.5
is retained rather than rounded away. Quality flags, duration, subgroup, and
matched-cohort analyses were treated as sensitivity analyses and are
reported in the supplementary evidence tables.

Figure~\ref{fig:mechanism} collects the principal mechanism and negative-control
results. The panels quantify label-predictive context, imprecise incremental
audio value, temporal feature-ranking instability, and source-target mismatch;
none alone identifies a causal shortcut encoded by the audio model.

\begin{figure}[!t]
\centering
\includegraphics[width=0.99\linewidth]{fig03_mechanism_and_incremental_value.pdf}
\caption{Mechanism and sensitivity evidence. (A) Metadata AUROC before and
after full shuffled-label retraining. (B) AUROC change after adding the
validation-ranked audio candidate to metadata baselines. (C) overlap of early
and late top-800 feature sets. (D) fraction of COUGHVID recordings inside the
source-domain classifier's 5th--95th percentile probability band.}
\label{fig:mechanism}
\end{figure}

\section{Discussion}

\subsection{Principal Findings}

The study yields five main findings. First, a multimodal Coswara model can
separate labels well in a participant-disjoint internal test, but that estimate
does not determine how a cough model will behave in another collection system.
Second, the external decline persisted across four conventional classifiers,
WavLM, and CNN--BiGRU. The consistency across representation and model families
makes it less plausible that the transfer gap is explained only by one weak
classifier. Third, feature selection was highly unstable across calendar
periods. Fourth, recording context and symptoms were strongly associated with
the label, and shuffled-label retraining removed that association. Fifth, on
the small aligned cohort, adding audio to metadata did not yield a precise or
statistically supported AUROC improvement.

These results do not imply that every internal result is invalid. They show
that the relevant claim must match the evaluation. AUROC 0.897 is evidence of
internal discrimination for one selected multimodal protocol. It is not an
estimate of multimodal performance on COUGHVID, because COUGHVID has no matched
breathing and speech inputs. Conversely, external AUROC near 0.5 is evidence of
poor portability under this source-target pair; it is not proof that no
disease-related acoustic information exists in any population.

\subsection{Relation to Prior Evidence}

The findings are consistent with the realistic-evaluation studies of Han
\emph{et al.} and Coppock \emph{et al.}, which showed that participant design,
demographic balance, symptoms, and enrolment change measured audio performance
\cite{han2022realistic,coppock2024audio}. Our study extends that concern within
a public multimodal Coswara workflow by adding forward and reverse calendar
tests, temporal feature-set overlap, explicit source-to-COUGHVID transfer,
deep-representation transfer, calibration, and target operating-point
analysis.

The comparison with the HST and DNDT/DNDF papers requires protocol discipline.
HST reported strong 10-fold results for dataset-specific Cambridge and
COUGHVID tasks \cite{aytekin2024hst}. The DNDT/DNDF paper reported Coswara and
COUGHVID within-dataset AUCs of 0.92 and 0.93, respectively, but its own
cross-dataset table reported Coswara-to-COUGHVID AUC around 0.53
\cite{islam2026robust}. Our conventional external range of 0.523--0.543 is
therefore not evidence that their internal classifiers were fabricated or
incorrect. It is evidence that within-dataset and cross-dataset numbers are not
interchangeable.

The low feature-set overlap complements the drift-adaptive motivation of
Ganitidis \emph{et al.} \cite{ganitidis2025drift}. Our reverse temporal result
also illustrates why AUROC alone is insufficient under a major prevalence or
case-mix change: ranking can remain high while precision, F1, and calibration
are poor. A deployment-oriented analysis must report both discrimination and
the consequences at a stated operating point.

\subsection{Implications for Model Development}

WavLM and CNN--BiGRU did not remove the external gap. Their results show that
the observed source--target difference was not confined to the conventional
feature bank. Architecture comparisons and transportability comparisons remain
separate questions because the cohort, split, threshold, and target population
also determine the estimate.

Fusion should likewise be interpreted as a fitted prediction layer rather than
an automatic benefit of having more modalities. Uniform averaging imposes
equal contribution; validation weighting uses observed validation quality; a
logistic stack learns coefficients from aligned validation predictions. The
test set cannot decide which fusion to report. This boundary is especially
important when participants have missing modalities or when one modality has a
different prevalence and quality distribution.

\subsection{Strengths}

The analysis preserves participant identities, prevents target data from
entering source model selection, distinguishes cough transfer from multimodal
fusion, reports calibration and base-rate-aware operating points, and uses
negative controls. External differences use independent-sample bootstrap
resampling, while DeLong tests are restricted to paired predictions. Results
are accompanied by split, quality, feature, and label-construction audits.

\subsection{Limitations}

The limitations define the scope of inference rather than invalidate the
study. First, Coswara and COUGHVID differ in recruitment, labels, prevalence,
devices, countries, and recording instructions. External failure therefore
cannot be attributed to one isolated mechanism. It is a transportability
result for the complete source-target shift.

Second, COUGHVID is cough-only. We can test whether a Coswara cough branch
transfers, but not whether cough--breath--speech fusion transfers. A second
external dataset with matched modalities and compatible labels would be needed
for that claim.

Third, the temporal protocols are affected by severe changes in label
prevalence and cohort composition. The reverse temporal AUROC is consequently
not a mirror-image estimate of forward deployment. It is retained as an
exploratory stress test with AUPRC and calibration, not presented as a superior
model.

Fourth, the metadata-plus-audio analysis has only 58--61 aligned test
participants across leading candidates. Its intervals are wide. The analysis
can rule out neither a modest benefit nor a modest harm, and it should not be
used as definitive proof that audio adds no information.

Fifth, metadata variables are measured imperfectly and cannot exhaust all
confounding paths. Strong metadata prediction and recording-year importance
establish association with collection context; they do not identify the exact
features used by an audio model. Matching depends on measured covariates and
available overlap and does not reproduce a randomized comparison.

Sixth, COUGHVID labels are imbalanced and may contain label noise. AUROC is
prevalence-invariant in definition but not immune to spectrum and label shift;
AUPRC, calibration, and precision are strongly prevalence-dependent. Both
discrimination and operational metrics are therefore reported.

Seventh, numerous exploratory analyses were performed after the primary model
development. They are labeled exploratory, and their isolated $p$-values are
not interpreted as confirmatory without multiplicity control. The primary
claims rest on the prespecified validation ladder and matched external model
families.

Eighth, the chronological evaluation reused a top-800 bank selected under the
original development split. It prevents later and external labels from entering
feature selection, but it is not a fully nested estimate of an end-to-end model
developed only in the early period. The temporal result is therefore a stress
test of the frozen workflow, not a prospective temporal-validation claim.

Finally, the study is retrospective and uses public crowdsourced audio. It does
not evaluate prospective workflow integration, patient outcomes, fairness in a
deployment population, or comparison with an approved clinical test. Decision
curves describe the present sample and utility assumptions; they do not show
that deployment would cause or prevent clinical harm.

\section{Conclusion}

A participant-disjoint multimodal model achieved strong internal
discrimination on Coswara, but matched cough models, WavLM, and CNN--BiGRU did
not retain that discrimination on COUGHVID. Calendar stress tests, unstable
feature selection, metadata predictiveness, calibration failure, and limited
source-target overlap provide complementary evidence that the benchmark is
sensitive to when, how, and from whom recordings were collected. These results
support protocol-matched reporting and external validation before respiratory-
audio performance is interpreted as portable screening evidence. They do not
establish a clinical diagnostic system or the absence of disease-related audio
information.


\appendices
\section{Reproducibility Details}

This appendix records the operational settings needed to interpret and
reproduce the completed analyses. It is intentionally more detailed than the
conference-length derivative of this manuscript.

\subsection{Cohort and Outcome Contract}

\begin{table}[!h]
\centering
\caption{Dataset roles and analysis units.}
\label{tab:data-contract}
\begin{tabularx}{\textwidth}{p{0.14\textwidth}p{0.18\textwidth}p{0.16\textwidth}X}
\toprule
Resource & Role and unit & Modalities & Outcome and boundary \\
\midrule
Coswara & Development and internal evaluation; participant is the split and
bootstrap unit & Cough, breathing, speech & Resolved positive/negative status;
unknown labels excluded from supervised analyses \\
COUGHVID & Frozen external target; recording UUID is the available unit &
Cough only & Processed \texttt{status\_SSL}; 285 positive and 8,046 negative
recordings; no target data used in source development \\
\bottomrule
\end{tabularx}
\end{table}

The external endpoint is not interchangeable with COUGHVID's expert audible-
pathology annotations or the dataset paper's cough-versus-noncough detector.
The source and target outcome fields also differ in ascertainment. External
performance therefore estimates transport to the complete COUGHVID analytic
setting, not a pure change in microphone acoustics.

\subsection{Signal and Feature Settings}

\begin{table}[!h]
\centering
\caption{Frozen conventional-audio processing settings.}
\label{tab:audio-settings}
\begin{tabular}{ll}
\toprule
Operation & Setting \\
\midrule
Waveform & Mono, 16~kHz, peak normalized \\
Low-energy trimming & 35-dB threshold \\
Active-region detector & 25-ms RMS frame, 10-ms hop, 0.25 peak ratio,
120-ms margin \\
Fixed duration & Cough 4~s; breathing 5~s; speech 5~s \\
Spectral grid & 400-sample FFT/window; 160-sample hop \\
Mel analysis & 64 bands, 20--8,000~Hz \\
Quality thresholds & Minimum 0.5~s cough or 1.0~s breath/speech; maximum
0.70 silence ratio and 0.01 clipping ratio \\
\bottomrule
\end{tabular}
\end{table}

The 10,140 numeric candidates comprised 6,377 ComParE-derived columns, 1,586
IS10-derived columns, and 2,177 project acoustic columns. The project bank
summarized MFCCs, mel energy, chroma, spectral contrast, tonal centroid,
zero-crossing rate, signal energy, spectral shape, and rhythm. These are
recording-level descriptions; they are not direct physiological measurements.

\subsection{Conventional Model Bank}

Every pipeline removed constant columns. A univariate ANOVA selector was then
fitted on training data using the model-specific retained percentage. Scaling
was used for logistic regression and SVC. SMOTE, where specified, used three
nearest neighbors and operated only after the training subset had been fixed.

\begingroup
\small
\begin{longtable}{L{0.22\textwidth}L{0.12\textwidth}L{0.56\textwidth}}
\caption{Principal conventional classifier settings.}
\label{tab:model-settings}\\
\toprule
Model & ANOVA features & Main settings \\
\midrule
\endfirsthead
\toprule
Model & ANOVA features & Main settings \\
\midrule
\endhead
L2 logistic & 80\% & $C=1$; class-balanced; maximum 4,000 iterations \\
RBF SVC & 60\% & $C=2$; scale-based $\gamma$; class-balanced; three-fold
sigmoid probability calibration \\
Extra Trees & 100\% & 700 trees and class balancing without SMOTE, or 800
trees with SMOTE; minimum leaf size 2 \\
Random Forest & 80\% & 700 trees with balanced subsampling, or 800 trees with
SMOTE; minimum leaf size 2 \\
XGBoost+SMOTE & 80\% & 800 trees; depth 3; learning rate 0.025; row and column
subsampling 0.85; minimum child weight 2; L2 penalty 2 \\
LightGBM+SMOTE & 80\% & 900 trees; 31 leaves; learning rate 0.025; minimum 20
child samples; row and column subsampling 0.85; L2 penalty 2 \\
CatBoost+SMOTE & 80\% & 900 iterations; depth 5; learning rate 0.025; log-loss
training objective \\
\bottomrule
\end{longtable}
\endgroup

The development bank also included SMOTE and non-SMOTE variants and a
validation-only Optuna search over logistic regression, Extra Trees, and
Random Forest. Model identity, modality combination, fusion method, and
threshold were selected from validation evidence. External labels did not
participate in this process.

\subsection{Deep Representation Settings}

WavLM Base Plus operated on 3-s windows with 50\% overlap and at most four
segments per recording. One augmented training copy was permitted. Training
used batch size 8, gradient accumulation 4, learning rates $2\times10^{-5}$
for the encoder and $10^{-4}$ for the prediction head, four unfrozen top
transformer layers, eight maximum epochs, and early stopping patience 3. The
CNN--BiGRU used log-mel spectrograms, batch size 16, learning rate $10^{-3}$,
and eight epochs. Both analyses selected checkpoints and thresholds using
Coswara validation data before frozen COUGHVID transfer.

\subsection{Interpretive Boundaries}

The following distinctions are part of the analysis contract:

\begin{itemize}
    \item the 0.897 row is internal cough--speech fusion, whereas COUGHVID tests
    only a cough branch;
    \item the early-to-late selected system is not the same model and modality
    combination as the internal selected fusion;
    \item independent source and target samples support independent bootstrap
    differences, not paired DeLong testing;
    \item metadata predictiveness is evidence of collection-context
    association, not proof that the audio model encoded the same variables;
    and
    \item post-hoc mechanism, subgroup, and utility analyses are exploratory.
\end{itemize}

\section{Protocol-Aware Comparison With Prior Work}

Table~\ref{tab:literature} is an estimand audit, not a classifier leaderboard.
An internal result trains and tests within one resource; a chronological result
changes calendar period; and a frozen external result applies a source-selected
model to a different resource. We use three comparison classes: A-close for a
substantially aligned endpoint, modality, and evaluation direction; B-partial
when at least one major protocol axis differs; and C-context when direct metric
ranking would be misleading. Even A-close rows are not formal head-to-head
tests because cohort releases and implementation details differ.

\begingroup
\scriptsize
\setlength{\tabcolsep}{3pt}
\begin{longtable}{L{0.15\textwidth}L{0.14\textwidth}L{0.17\textwidth}L{0.12\textwidth}L{0.15\textwidth}L{0.17\textwidth}}
\caption{Full-text-audited respiratory-audio literature, organized by the
quantity each experiment estimates. Target use includes fitting, cohort
selection, representation selection, thresholding, or calibration.}
\label{tab:literature}\\
\toprule
Study & Data and inputs & Unit and protocol & Reported result & Target-data use & Comparison boundary \\
\midrule
\endfirsthead
\toprule
Study & Data and inputs & Unit and protocol & Reported result & Target-data use & Comparison boundary \\
\midrule
\endhead
Bhattacharya et al., \emph{Coswara} \cite{bhattacharya2023coswara} & Coswara; nine audio tasks plus symptoms & Participant; 70/15/15 internal split & AUROC 0.915 (95\% CI 0.885--0.941) & Internal validation and test & B-partial: different snapshot, tasks, symptoms, and model \\
Orlandic et al., \emph{COUGHVID} \cite{orlandic2021coughvid} & COUGHVID cough events & Recording; nested 10-fold internal CV & AUC $0.964\pm0.033$ & Internal & C-context: cough-versus-noncough detection, not COVID classification \\
Xia et al., \emph{COVID-19 Sounds} \cite{xia2021covid} & Cough, breath, and voice & Participant; speaker-independent 7:1:2 split & Multimodal AUROC 0.71 (95\% CI 0.65--0.76) & Internal & B-partial: multimodal but a different cohort and label process \\
Chetupalli et al. \cite{chetupalli2023multimodal} & Coswara deep breath, heavy cough, counting, symptoms & Participant; subject-disjoint 80/20 with development CV & Audio-only 0.88; audio+symptoms 0.963 & Internal & A-close internal context for 0.88; symptom-assisted 0.963 is C-context \\
Truong et al., FAIR \cite{truong2024fair} & Coswara; seven sound instances & Participant; fixed test plus rotating development folds and two seeds & AUROC $0.8658\pm0.0115$ & Internal & A-close internal context; different cohort and seven-input architecture \\
Han et al. \cite{han2022realistic} & COVID-19 Sounds; cough, breath, voice & Participant; realistic and deliberately biased internal designs & Realistic multimodal 0.71 (95\% CI 0.65--0.77); biased designs up to 0.90 & Internal & B-partial: direct evidence that cohort construction changes AUROC \\
Coppock et al. \cite{coppock2024audio} & 67,842 PCR-linked UK participants; cough, exhalation, sentence & Participant; random, matched, longitudinal, and longitudinal-matched & Sentence 0.846 random versus 0.619 matched & Internal, prespecified covariate controls & B-partial: strongest methodological comparator; substantially larger PCR cohort \\
Ganitidis et al. \cite{ganitidis2025drift} & Coswara and COVID-19 Sounds cough & Participant-separated chronological 70/30 plus adaptation experiments & Coswara 0.668 to 0.597; COVID-19 Sounds 0.691 to 0.607 & Later labels used for evaluation and adaptation & B-partial: close temporal question, but cough-only and adaptation-focused \\
Kim and Lee \cite{kim2024longitudinal} & Cambridge source; Virufy and Coswara periods & Cough; dataset and period transfer; participant grouping not explicit & Cambridge 0.9346; Coswara 0.8250, 0.7724, 0.5509 by period & No target retraining reported & B-partial: acquisition, calendar, presumed variant, and prevalence change together \\
Grant et al. \cite{grant2022deployment} & DiCOVA/Coswara source; device and source targets & Recording/participant as available; source CV and stress tests & Web cough 0.42 versus Android cough 0.75; mean 0.77 & Frozen and target-optimized thresholds both reported & B-partial: deployment-mismatch evidence; different source-target pairs \\
Atmaja et al. \cite{atmaja2025crossdataset} & Pooled Coswara and COUGHVID training; ComParE test & Cough segments; fixed development and 154-segment test & Unweighted accuracy 0.8819 & The same test supported choices of split, segmentation, augmentation, and mixup & C-context numerically: pooled-source training and test-guided configuration; result is not AUROC \\
Aytekin et al., HST \cite{aytekin2024hst} & Cambridge tasks and separate balanced COUGHVID task & Dataset-specific repeated participant-disjoint internal partitions & COUGHVID AUROC $0.90\pm0.01$ & COUGHVID used for training and selection & C-context for transfer; valid internal architecture comparison only \\
Islam et al., DNDT/DNDF \cite{islam2026robust} & Five cough datasets; selected balanced cohorts & Stratified internal CV and explicit train-one/test-another transfer & Internal: Coswara 0.92, COUGHVID 0.93; cross-data: 0.53 & Target cohort balanced by label; no target fitting in transfer & A-close for Coswara-to-COUGHVID direction and cough modality; target is 680/680 and participant grouping is unreported \\
Pahar et al. \cite{pahar2022transfer} & Coswara; Sarcos target; separate cough, breath, speech & Subject-aggregated nested leave-$p$-out; external target has 44 subjects & Coswara cough 0.982; Coswara$\rightarrow$Sarcos 0.954 & No target fitting reported & B-partial: genuine transfer, but different and very small target \\
Avila et al. \cite{avila2021feature} & DiCOVA/Coswara cough & Fixed participant folds and 233-recording blind test & Blind fusion 0.808; nonnested feature CV 0.96 fell to 0.768 when nested & Blind test excluded from selection & B-partial and direct precedent for feature-selection optimism \\
Laguarta et al. \cite{laguarta2020cough} & Private 5,320-subject forced-cough cohort & Subject; random 80/20 internal split & AUROC 0.97 on the official-test subset & Internal & C-context: private internal task with heterogeneous label ascertainment \\
Chowdhury et al. \cite{chowdhury2022mcdm} & Several cough datasets, separate or merged & Stratified 10-fold CV; participant grouping unreported & Coswara about 0.64--0.66; other internal rows up to 0.97--0.98 & Internal or pooled & C-context for transfer; no frozen train-one/test-another estimate \\
Celik, \emph{CovidCoughNet} \cite{celik2023covidcoughnet} & COUGHVID and Coswara; modalities modeled separately & Recording-level 80/20 split; participant and augmentation-parent grouping unreported & COUGHVID binary AUC 0.9505; Coswara cough 0.9848 & Test set used across pitch-shift comparisons & C-context: merged symptomatic/COVID endpoint, no multimodal fusion or external transfer \\
Hussain et al., \emph{Cough2COVID-19} \cite{hussain2024cough2covid} & Coswara, Virufy, ComParE development; selected COUGHVID target & Reported multi-source leave-one-dataset-out cough evaluation & COUGHVID AUROC 0.981 & Labeled results from all four resources informed representation-family ranking & B-partial: important counterexample; multi-source training, selected 651/660 target, and no uncertainty \\
de Brito et al. \cite{debrito2025audiomae} & Coswara and COUGHVID cough & Stratified 80/20 and bidirectional transfer after demographic balancing & Cross-data AUROC 0.51--0.68 across Audio-MAE/PANN & Target labels define balanced evaluation cohorts & A-close scientifically for Coswara-to-COUGHVID transfer, but non-peer-reviewed and not the full target distribution \\
Lin et al., \emph{SympCoughNet} \cite{lin2025sympcoughnet} & UK cough plus 14 symptoms & Sample; stratified 70/15/15; participant separation not reported & Symptom+audio 0.9474; audio-only ablation 0.8127 & Internal & C-context for headline; B-partial for audio-only ablation \\
\midrule
Present study & Coswara cough, breath, speech; full known-label COUGHVID cough target & Source participant; target recording; internal, chronological, and frozen transfer & Internal multimodal 0.897; matched cough external 0.523--0.543 & No target labels in source model or threshold; separate recalibration sensitivity disclosed & Adds common audit controls, but does not externally validate multimodal fusion \\
\bottomrule
\end{longtable}
\endgroup

The closest internal multimodal context is Chetupalli's audio-only fusion at
0.88 and FAIR at $0.8658\pm0.0115$; the present selected internal estimate is
0.897. These values are contextual, not superiority tests. The closest
peer-reviewed frozen direction is DNDT/DNDF's Coswara-to-COUGHVID AUROC of
0.53, consistent with the present conventional range of 0.523--0.543. The
Audio-MAE/PANN technical report independently places deep transfer in the
0.51--0.68 range, while the present WavLM and CNN--BiGRU estimates are 0.484
and 0.548.

Cough2COVID-19 is a material counterexample rather than a result to dismiss. It
reports 0.981 for multi-source development followed by a selected balanced
COUGHVID target. That estimand differs from single-source Coswara transfer to
all 8,331 known-label COUGHVID recordings, and labeled results from all four
resources informed its representation-family ranking. The contrast therefore
motivates a future controlled factorial study of source diversity, target
selection, and representation choice; it does not establish that either
reported estimate is erroneous.


\section*{Data and Code Availability}
The analysis uses the public Coswara and COUGHVID resources under their
respective access and use conditions. The executable analysis code, frozen
configuration, split audits, and aggregate evidence tables will be released
with the accepted version. Raw recordings are not redistributed by this work.

\section*{Ethics Statement}
This study is a secondary analysis of previously collected datasets. The
original dataset publications describe their participant information,
consent, and ethics procedures. No new participants were recruited and no new
identifying information was collected. The exact wording and identifiers from
the source publications will be reproduced in the submission version after
author verification.

\bibliographystyle{IEEEtran}
\bibliography{references}

\end{document}
